%Chargement des paquets
\usepackage{amsmath}
\usepackage{amsfonts}
\usepackage{amssymb}
\usepackage{enumerate}
\usepackage{mathtools}
\usepackage{bbm}
\usepackage{xparse, etoolbox}
\usepackage{enumerate}
\usepackage{enumitem}
\usepackage{mathabx}
\usepackage{babel}[french]

%environment
\usepackage{amsthm}
\usepackage{keytheorems}
\usepackage{hyperref} 

%Interface théorème
\newkeytheorem{lem}[
	name = Lemme
]

\newkeytheorem{prop}[
	name = Proposition
]

\newkeytheorem{thm}[
	name = Théorème
]

\newkeytheorem{coro}[
	name = Corollaire
]

\renewcommand*{\proofname}{Démonstration}

\theoremstyle{definition}
\newtheorem{defn}{Définition}

\theoremstyle{definition}
\newtheorem*{nota}{Notation}

\theoremstyle{definition}
\newtheorem*{limi}{Liminaire}

\theoremstyle{remark}
\newtheorem{rmq}{Remarque}

\theoremstyle{plain}
\newtheorem*{fait}{Fait}


%Ensembles de nombres
\newcommand{\N} {\mathbb{N}}
\newcommand{\Ne}{\N^\ast}
\newcommand{\Z} {\mathbb{Z}}
\newcommand{\Q} {\mathbb{Q}}
\newcommand{\R} {\mathbb{R}}
\newcommand{\Rb}{\overline{\mathbb{R}}}
\newcommand{\Rp}{\R_+}
\newcommand{\Rm}{\R_-}
\newcommand{\K} {\mathbb{K}}
\newcommand{\Ket} {\mathbb{K^{\ast}}}
\newcommand{\Cx}{\mathbb{C}}

%Ensembles, tribus
\newcommand{\A}{\mathbb{A}}
\renewcommand{\P}{\mathcal{P}}
\newcommand{\T}{\mathcal{T}}
\newcommand{\F}{\mathcal{F}}
\newcommand{\C} {\mathcal{C}}
\newcommand{\ouv}{\mathcal{O}}
\newcommand{\B}{\mathcal{B}}
\newcommand{\M}{\mathcal{M}}
\newcommand{\etag}{\mathcal{E}}
\newcommand{\etagOp}{\etag^{+}(\Omega)}
\newcommand{\etagO}{\etag(\Omega)}
%%Lp avant quotientage
\newcommand{\Lic}{\mathcal{L}^1}
\newcommand{\Lpc}{\mathcal{L}^p}
\newcommand{\Linfc}{\mathcal{L}^{\infty}}
\newcommand{\Lifull}{\Lic(\Omega, \B, m)}
\newcommand{\Lpfull}{\Lpc(\Omega, \B, m)}
\newcommand{\Linffull}{\Linfc(\Omega, \B, m)}
%%Lp après quotientage
\newcommand{\Li}{L^1}
\newcommand{\Ld}{L^2}
\newcommand{\Lp}{L^p}
\newcommand{\Linf}{L^{\infty}}
\newcommand{\Listd}{\Li(\Omega, m)} 
\newcommand{\Ldstd}{\Ld(\Omega, m)} 
\newcommand{\Lpstd}{\Lp(\Omega, m)} 

%Quelques opérateurs
\DeclareMathOperator{\codim}{codim}
\DeclareMathOperator{\rg}{rg}
\DeclareMathOperator{\tr}{tr}
\DeclareMathOperator{\vect}{Vect}
\DeclareMathOperator{\ima}{Im}
\DeclareMathOperator{\id}{Id}
\DeclareMathOperator{\sign}{sign}
\DeclareMathOperator{\ext}{ext}
\newcommand{\ind}{\mathbbm{1}}
\newcommand*{\spr}[2]{\left(\,#1\,\middle|\,#2\,\right)}
\newcommand{\interior}[1]{%
  {\kern0pt#1}^{\mathrm{o}}%
}
\DeclareMathOperator{\card}{card}
\DeclareMathOperator{\ppcm}{ppcm}

%Commandes de pure flemme
\newcommand{\veps}{\varepsilon}
\newcommand{\intab}{\int_{a}^{b}}
\newcommand{\intR}{\int_{-\infty}^{\infty}}
\newcommand{\intinf}{\int_{- \infty}^{+ \infty}}
\newcommand{\equi}{\Leftrightarrow}
\newcommand{\Kev}{$\K$-espace vectoriel }
\newcommand{\Kevs}{$\K$-espaces vectoriels }
\newcommand{\Rev}{$\R$-espace vectoriel }
\newcommand{\Revs}{$\R$-espaces vectoriels }
\newcommand{\Cev}{$\Cx$-espace vectoriel }
\newcommand{\Cevs}{$\Cx$-espaces vectoriels }
\newcommand{\evn}{(E, \norm{.})}
\newcommand{\interf}[2]{[#1,#2]}
\newcommand{\limb}{\lim\limits_}
\newcommand{\lebn}{\lambda_n}
\newcommand{\pp}{\text{~presque partout}}
\newcommand{\derivp}[2]{\frac{\partial #1}{\partial #2}}
\newcommand{\famiu}{(u_i)_{i \in I}}
\newcommand{\sev}{\text{~s.e.v.}}
\newcommand{\Mnp}{M_{n,p}(\K)}
\newcommand{\Mn}{M_{n}(\K)}
\newcommand{\tq}{\text{~t.q.~}}
\newcommand{\mq}{~montrer~que~}
\newcommand{\et}{\quad \text{et} \quad}
\newcommand{\ou}{\text{~ou~}}
\newcommand{\Kx}{\K[X]}


%Commandes de gars chelou
\renewcommand{\subset}{\subseteq}

%Commande restriction, ty duckduckgo
\def\restr#1#2{\mathchoice
              {\setbox1\hbox{${\displaystyle #1}_{\scriptstyle #2}$}
              \restraux{#1}{#2}}
              {\setbox1\hbox{${\textstyle #1}_{\scriptstyle #2}$}
              \restraux{#1}{#2}}
              {\setbox1\hbox{${\scriptstyle #1}_{\scriptscriptstyle #2}$}
              \restraux{#1}{#2}}
              {\setbox1\hbox{${\scriptscriptstyle #1}_{\scriptscriptstyle #2}$}
              \restraux{#1}{#2}}}
\def\restraux#1#2{{#1\,\smash{\vrule height .8\ht1 depth .85\dp1}}_{\,#2}} 

%Quelques normes
\DeclarePairedDelimiter\abs{\lvert}{\rvert}
\DeclarePairedDelimiter\labs{\bigl\lvert}{\bigr\rvert}
\DeclarePairedDelimiter\norm{\lVert}{\rVert}
\DeclarePairedDelimiterXPP\normi[1]{}\lVert\rVert{_1}{\ifblank{#1}{\:\cdot\:}{#1}}
\DeclarePairedDelimiterXPP\normd[1]{}\lVert\rVert{_2}{\ifblank{#1}{\:\cdot\:}{#1}}
\DeclarePairedDelimiterXPP\normp[1]{}\lVert\rVert{_p}{\ifblank{#1}{\:\cdot\:}{#1}}
\DeclarePairedDelimiterXPP\normq[1]{}\lVert\rVert{_q}{\ifblank{#1}{\:\cdot\:}{#1}}
\DeclarePairedDelimiterXPP\norminf[1]{}\lVert\rVert{_\infty}{\ifblank{#1}{\:\cdot\:}{#1}}

%Bracket en angle
\DeclarePairedDelimiter\angl{\langle}{\rangle}

%Set, parenthèses
\newcommand{\set}[1]{\{ #1 \}}
\newcommand{\bigset}[1]{\bigl\{ #1 \bigr\}}
\newcommand{\Bigset}[1]{\Bigl\{ #1 \Bigr\}}
\newcommand{\bigpar}[1]{\bigl( #1 \bigr)}
\newcommand{\Bigpar}[1]{\Bigl( #1 \Bigr)}

%Opitmisation des opérateurs de comparaison
\DeclarePairedDelimiter\tio{o (}{)}
\DeclarePairedDelimiter\gro{O (}{)}


\begin{document}

\pagestyle{fancy}
\fancyhead[C]{\textbf{La méthode de Newton}}

\underline{Objectif :} recherche d'une racine d'une fonction $\C^2$ sur un compact. \\

\underline{Heuristique :} la méthode de Newton consiste à prendre un point $x_0$ sur le graphe de la fonction $f$ 
puis à remplacer la fonction par sa tangente en ce point, c'est-à-dire par le droite d'équation $y = f'(x_0)(x-x_0) + f(x_0)$.
Si $f$ n'est pas localement constante, la tangente coupe l'axe des abscisses en un nouveau point $x_1 = x_0 + \dfrac{f(x_0)}{f'(x_0)}$ 
ce qui permet, à priori, de créer une suite. \\
Lorsque la fonction $f$ est strictement croissante, toutes ses tangents vont effectivement couper l'axe des abscisses ce qui permet de 
définir effectivement une suite $(x_n)_\N$. \\
Si de plus $f$ est strictement convexe (c'est-à-dire que son graphe est au-dessus de toutes ses tangentes), alors on peut garantir que 
tous les points de la suite sont plus grand que $a$, et même, en conditionnant le choix de $x_0$ de garantir que la suite converge vers le
zéro de la fonction. 

\begin{thm}
Soit $f:[c,d] \to \R$ une fonction de classe $\C^2$ ; on suppose que $c < d, f(c) < 0 < f(d)$ et f$'(x) >0$ pour tout $x \in [c,d]$.  \\
On considère la suite $(x_n)$, définie sur $\N$ par :
\[ x_0 \in [c,d] \et x_{n+1} = F(x_n), \text{ avec } F(x) = x - \dfrac{f(x)}{f'(x)} . \]
Alors, $f$ admet un unique zéro sur $[c,d]$, que l'on note $a$, et :
\begin{enumerate}[label=(\roman*)]
\item il existe un intervalle $I$ centré en $a$, stable par $F$, et tel que, pour tout $x_0 \in I$, la suite $(x_n)_\N$ converge
quadratiquement vers $a$.
\item si de plus $f''(x)>0$ pour tout $x \in [c,d]$, alors pour tout $x_0 \in [a,d]$, la suite $(x_n)_\N$ est strictement décroissante 
(ou constante) et converge quadratiquement vers $a$.
\end{enumerate}
\end{thm}

\begin{proof}
$f$ étant de classe $\C^2$ sur $[c,d]$, en particulier continue et strictement continue sur $[c,d]$, elle admet bien un zéro unique, $a$, 
sur cet intervalle. 

\begin{enumerate}[label=(\roman*)]

\item En utilisant la formule de Taylor sur l'intervalle d'extrémités $x$ et $a$, 
on montre qu'il existe $z$ dans l'intérieur de cet intervalle tel que :
\begin{gather*}
f(a) = f(x) + f'(x)(a-x) + \dfrac12 f''(z)(a-x)^2 = 0 \\
\text{soit, } f'(x)(x-a) - f(x) = \dfrac12 f''(z)(x-a)^2
\end{gather*}
on en déduit que : 
\begin{align*}
F(x) - a & = (x-a) - \dfrac{f(x)}{f'(x)} \\
 & = \dfrac{f'(x)(x-a) - f(x)}{f'(x)} \\
 & = \dfrac12 \dfrac{f''(z)}{f'(x)}(x-a)^2 \qquad (*)
\end{align*}

La fonction $x \mapsto \dfrac12 \dfrac{f''(x)}{f'(x)}$ étant continue sur $[c,d]$ elle est bornée, notons $C$ sa borne (en valeur absolue), 
on en déduit que :
\[ \abs{F(x)-a} \leq C \abs{x-a}^2 \qquad (**) \]
Ainsi, en prenant $\alpha > 0$ tel que $C \alpha < 1$ et $I = [a-\alpha ; a+\alpha] \subset [c,d]$, on obtient, pour tout $x \in I$ :
\[ \abs{F(x) - a} < C\abs{x-a}^2 < C \alpha^2 < \alpha , \]
ce qui démontre que l'intervalle $I$ est stable par $F$. 
De plus, pour $x_0 \in I$, comme $x_{n+1} = F(x_n)$ on montre par récurrence que :
\[ C \abs{x_n - a}  \leq  ( C \alpha)^{2n} , \]
ce qui garantit la convergence quadratique car $C \alpha < 1$.

\item Si de plus $f''>0$ sur $[c,d]$, on déduit de (*) que $\forall x \in [c,d], F(x) > a$.
Par ailleurs, $F(x)-x = - \dfrac{f(x)}{f'(x)} \leq 0$ avec égalité si, et seulement si, $x=a$. 
On en déduit que l'intervalle $[a,d]$ est stable par $F$ et que, si $x_0 = a$ alors la suite $(x_n)_\N$ est constante, égale à $a$, 
sinon, pour $x_0 \in ]a,d]$, la suite est strictement décroissante. \\
Ainsi, la suite converge sur $[a,d]$, et sa limite est un point fixe de $F$ sur cet intervalle soit $a$. \\
De plus  l'inégalité (**) étant toujours valable, on en déduit que la convergence est bien quadratique.
\end{enumerate}
\end{proof}


\end{document}